% AA-25 | Iteracja XIII; szczegolowy audyt historyczny 2026-11-20, partia 02.
\chapter{90. twierdzenie Hilberta i równania normowe}
\label{aa:ch:hilbert90}

\section{Motywacja: kiedy element jest ilorazem swoich sprzężeń?}
Niech $L/K$ będzie skończonym rozszerzeniem Galois, a $G=\operatorname{Gal}(L/K)$
jego grupą Galois. Jeżeli $b\in L^\times$ i $\sigma\in G$, to element
$a=b/\sigma(b)$ ma szczególną własność: gdy $G=\langle\sigma\rangle$, jego
norma z $L$ do $K$ wynosi $1$. W istocie kolejne czynniki
$a,\sigma(a),\ldots,\sigma^{n-1}(a)$ znoszą się po wymnożeniu.
Pojawia się pytanie odwrotne: czy warunek $N_{L/K}(a)=1$ wystarcza,
aby znaleźć takie $b$? Odpowiedź daje \emph{90. twierdzenie Hilberta}.
Wykażemy też wersję addytywną, w której normę zastępuje ślad,
a iloraz $b/\sigma(b)$ różnica $b-\sigma(b)$.

Normę $N_{L/K}$ i ślad $\operatorname{Tr}_{L/K}$ definiowaliśmy
w rozdziale~\ref{aa:ch:trace-norm}. Przypomnijmy, że są one określone
za pomocą wyznacznika i śladu $K$-liniowego odwzorowania
$L\ni x\mapsto ax\in L$; nie zakładamy przy tym, że
$\operatorname{char}K=0$. Dla automorfizmu $\sigma$ symbol $\sigma^j$
oznacza $j$-krotne złożenie, a $\sigma^0=\operatorname{id}_L$.

\section{Niezależność automorfizmów i przygotowanie}
Twierdzenie multiplikatywne będzie wymagało znalezienia elementu,
na którym pewna suma automorfizmów nie znika. Nie można tego
uzasadnić samym faktem, że każdy jej składnik jest niezerowy:
suma niezerowych odwzorowań bywa zerowa. Potrzebujemy silniejszego lematu.

\begin{lemat}[Niezależność różnych homomorfizmów ciał]\label{aa:lem:dedekind-independent}
Niech $F,E$ będą ciałami, $r\geq1$ liczbą całkowitą,
a $\sigma_1,\ldots,\sigma_r:F\to E$ parami różnymi homomorfizmami ciał.
Dla dowolnych $c_1,\ldots,c_r\in E$, jeżeli
$\sum_{i=1}^r c_i\sigma_i(x)=0$ dla każdego $x\in F$,
to $c_1=\cdots=c_r=0$.
\end{lemat}
\begin{proof}
Załóżmy wbrew tezie, że istnieje niezerowy układ współczynników
spełniający wskazaną zależność. Spośród wszystkich takich zależności
wybierzmy tę o najmniejszej liczbie niezerowych współczynników.
Po usunięciu składników zerowych i zmianie numeracji możemy zapisać ją jako
\[
\sum_{i=1}^{s}c_i\sigma_i(x)=0\qquad(x\in F),
\quad c_i\neq0\ (1\leq i\leq s).
\]
Musi być $s\geq2$: dla $s=1$ podstawienie $x=1$ daje
$c_1\sigma_1(1)=c_1=0$, sprzecznie z wyborem.
Homomorfizmy $\sigma_1$ i $\sigma_s$ są różne, zatem
istnieje $u\in F$ spełniające $\sigma_1(u)\neq\sigma_s(u)$.

Dla każdego $x\in F$ zastosujmy początkową zależność do $ux$.
Ponieważ każdy $\sigma_i$ zachowuje mnożenie, otrzymujemy
$\sum_{i=1}^{s}c_i\sigma_i(u)\sigma_i(x)=0$.
Pomnóżmy osobno początkową zależność dla $x$ przez
$\sigma_s(u)\in E$ i odejmijmy ją od otrzymanej równości.
Dostajemy dla wszystkich $x\in F$ zależność
\[
\sum_{i=1}^{s}c_i\bigl(\sigma_i(u)-\sigma_s(u)\bigr)\sigma_i(x)=0.
\]
Współczynnik przy $\sigma_s$ jest zerowy, natomiast przy
$\sigma_1$ wynosi $c_1(\sigma_1(u)-\sigma_s(u))\neq0$,
ponieważ $E$ jest ciałem. Powstała zatem niezerowa zależność
z co najwyżej $s-1$ niezerowymi współczynnikami. Przeczy to minimalności
$s$ i kończy dowód.
\end{proof}

\begin{uwaga}[Zakres lematu]
Homomorfizmy $\sigma_i$ nie muszą być surjektywne ani automorfizmami
tego samego ciała. W dalszych zastosowaniach przyjmiemy $F=E=L$
i wybierzemy $\sigma_i$ spośród elementów grupy Galois.
Wtedy suma $\sum_i c_i\sigma_i$ jest odwzorowaniem $L\to L$,
a lemat stwierdza, że nie może być identycznie zerowa,
jeżeli co najmniej jeden współczynnik $c_i$ jest niezerowy.
\end{uwaga}

\begin{lemat}[Norma i ślad w rozszerzeniu Galois]\label{aa:lem:galois-norm-product}
Niech $L/K$ będzie skończonym rozszerzeniem Galois oraz
$G=\operatorname{Gal}(L/K)$. Wówczas dla każdego $a\in L$ zachodzą
\[
N_{L/K}(a)=\prod_{\tau\in G}\tau(a),\qquad
\operatorname{Tr}_{L/K}(a)=\sum_{\tau\in G}\tau(a).
\]
\end{lemat}
\begin{proof}
Niech $n=[L:K]$. Ponieważ rozszerzenie $L/K$ jest rozdzielcze,
ma dokładnie $n$ różnych $K$-zanurzeń w ustalone domknięcie
algebraiczne $\overline K$. Normalność oznacza, że obraz każdego
z tych zanurzeń jest zawarty w $L$. Każde takie zanurzenie jest
$K$-liniowe i iniektywne; jako iniektywne odwzorowanie skończenie
wymiarowej przestrzeni $L$ w siebie jest również surjektywne,
a zatem jest elementem $G$.
Stąd $|G|=n$. Wzory na normę i ślad jako iloczyn i sumę
obrazów elementu przy wszystkich $K$-zanurzeniach,
udowodnione w rozdziale~\ref{aa:ch:trace-norm},
dają obie żądane równości. Wzór iloczynowy obejmuje również $a=0$.
\end{proof}

\section{Multiplikatywna postać 90. twierdzenia Hilberta}
\begin{twierdzenie}[90. twierdzenie Hilberta — postać multiplikatywna]\label{aa:tw:hilbert90-multiplicative}
Niech $L/K$ będzie skończonym cyklicznym rozszerzeniem Galois,
$G=\langle\sigma\rangle$ grupą rzędu $n\geq1$, a $a\in L^\times$.
Następujące warunki są równoważne:
\begin{enumerate}[label=\textup{(\roman*)}]
\item $N_{L/K}(a)=1$;
\item istnieje $b\in L^\times$ takie, że $a=b/\sigma(b)$.
\end{enumerate}
\end{twierdzenie}
\begin{proof}
\emph{Kierunek (ii)$\Rightarrow$(i).}
Niech $a=b/\sigma(b)$, gdzie $b\neq0$. Każdy automorfizm
ciała zachowuje niezerowość, więc wszystkie mianowniki
$\sigma^j(b)$ są niezerowe. Korzystając z lematu
\ref{aa:lem:galois-norm-product} oraz $\sigma^n=\operatorname{id}_L$,
otrzymujemy
\[
N_{L/K}(a)=\prod_{j=0}^{n-1}\frac{\sigma^j(b)}{\sigma^{j+1}(b)}
=\frac{b}{\sigma^n(b)}=1.
\]

\emph{Kierunek (i)$\Rightarrow$(ii).}
Załóżmy, że $N_{L/K}(a)=1$. Zdefiniujmy
$ p_0=1$ oraz, dla $1\leq j\leq n$,
$ p_j=\prod_{i=0}^{j-1}\sigma^i(a)$.
Ponieważ $a\neq0$, wszystkie $p_j$ są niezerowe;
ponadto z założenia i lematu o normie wynika $p_n=1$.
Dla $0\leq j<n$ obliczamy, osobno wypisując przesunięcie indeksu,
\[
\sigma(p_j)=\prod_{i=0}^{j-1}\sigma^{i+1}(a)
=\frac{p_{j+1}}{a};
\]
przy $j=0$ obie strony są równe $1$. Definiujemy odwzorowanie
$B:L\to L$ wzorem
\[
B(x)=\sum_{j=0}^{n-1}p_j\sigma^j(x),\qquad x\in L.
\]
Automorfizmy $\operatorname{id},\sigma,\ldots,\sigma^{n-1}$
są parami różne. Z lematu~\ref{aa:lem:dedekind-independent}
i faktu, że $p_0=1\neq0$, wynika, że $B$ nie jest odwzorowaniem
identycznie zerowym. Istnieje więc $x\in L$ takie, że
$b:=B(x)\neq0$; tutaj wskazaliśmy zarówno sposób zdefiniowania
$b$, jak i przyczynę jego niezerowości.

Zastosujmy $\sigma$ do każdego składnika $b$ i skorzystajmy
z tożsamości dotyczącej $\sigma(p_j)$:
\[
\begin{aligned}
a\sigma(b)
&=a\sum_{j=0}^{n-1}\sigma(p_j)\sigma^{j+1}(x)
 =\sum_{j=0}^{n-1}p_{j+1}\sigma^{j+1}(x)\\
&=\sum_{j=1}^{n-1}p_j\sigma^j(x)
   +p_n\sigma^n(x)
 =\sum_{j=1}^{n-1}p_j\sigma^j(x)+x=B(x)=b.
\end{aligned}
\]
Ostatnie przejście używa $p_n=1$, $\sigma^n(x)=x$ i $p_0=1$.
Ponieważ $b\neq0$, także $\sigma(b)\neq0$, więc wolno podzielić
przez $\sigma(b)$. Dostajemy $a=b/\sigma(b)$.
Argument obejmuje również $n=1$: wtedy $G$ jest trywialna,
$N_{L/K}(a)=a=1$, a można wybrać $b=1$.
\end{proof}

\begin{przyklad}[Norma zespolona]\label{aa:prz:hilbert90-complex}
Rozpatrujemy rozszerzenie $\mathbb C/\mathbb R$ stopnia $2$;
jego grupę Galois generuje sprzężenie $\sigma(z)=\overline z$.
Dla $z=x+iy$, $x,y\in\mathbb R$, mamy
$N_{\mathbb C/\mathbb R}(z)=(x+iy)(x-iy)=x^2+y^2=|z|^2$.
Warunek normy $1$ jest zatem równoważny $|a|=1$.
Jeśli $a\neq-1$, wybierzmy $b=1+a$. Jest on niezerowy, a ponieważ
$a\overline a=1$, zachodzi $\overline a=a^{-1}$ oraz
\[
\frac{b}{\overline b}
=\frac{1+a}{1+\overline a}
=\frac{1+a}{1+a^{-1}}=a.
\]
Dla $a=-1$ wyrażenie $1+a$ znika, więc trzeba zmienić wybór:
$b=i$ daje $b/\overline b=i/(-i)=-1$.
Przykład pokazuje również różnicę terminologiczną: norma
rozszerzenia $\mathbb C/\mathbb R$ nie jest modułem $z$, lecz
jego kwadratem. Dla $z=2+3i$ norma wynosi $13$, a moduł $\sqrt{13}$.
\end{przyklad}

\begin{przyklad}[Norma w rozszerzeniu kwadratowym]\label{aa:prz:hilbert90-quadratic}
Niech $K$ będzie ciałem charakterystyki różnej od $2$, a
$d\in K^\times$ nie będzie kwadratem w $K$. Wtedy
$L=K(\sqrt d)$ ma bazę $(1,\sqrt d)$ nad $K$, a
$\sigma(u+v\sqrt d)=u-v\sqrt d$. Dla $u,v\in K$ obliczamy
\[
N_{L/K}(u+v\sqrt d)=(u+v\sqrt d)(u-v\sqrt d)=u^2-dv^2.
\]
Jeśli $N(a)=1$ i $a\neq-1$, można \emph{jawnie} wziąć $b=1+a$,
ponieważ $\sigma(a)=a^{-1}$ i ten sam rachunek co w przykładzie
zespolonym daje $b/\sigma(b)=a$. Jeśli $a=-1$, wybieramy
$b=\sqrt d$, dla którego $\sigma(b)=-b$.

Na przykład dla $K=\mathbb Q$, $d=2$ i $a=3+2\sqrt2$ mamy
$N(a)=9-8=1$. Bierzemy $b=1+a=4+2\sqrt2$ i sprawdzamy:
\[
\frac{4+2\sqrt2}{4-2\sqrt2}
=\frac{(4+2\sqrt2)^2}{16-8}
=\frac{24+16\sqrt2}{8}=3+2\sqrt2.
\]
Ważne jest rozróżnienie dwóch pytań: 90. twierdzenie Hilberta
opisuje \emph{jądro} normy, natomiast równanie $N(b)=c$
pyta, czy dane $c\in K^\times$ należy do \emph{obrazu} normy.
Nie wynika stąd, że dowolne $c$ jest normą.
Dla $\mathbb C/\mathbb R$ mamy
$N_{\mathbb C/\mathbb R}(\mathbb C^\times)=\mathbb R_{>0}$,
ponieważ $x^2+y^2>0$ dla $(x,y)\neq(0,0)$,
a każdą dodatnią liczbę $c$ otrzymujemy jako normę $\sqrt c$.
W szczególności $-1$ nie jest taką normą.
\end{przyklad}

\begin{przyklad}[Multiplikatywne równanie w ciele czteroelementowym]\label{aa:prz:hilbert90-f4}
Niech $K=\mathbb F_2$, a $L=\mathbb F_2[\theta]$, gdzie
$\theta^2+\theta+1=0$. Wielomian $X^2+X+1$ nie ma
pierwiastka w $K$, zatem $L/K$ jest rozszerzeniem Galois
stopnia $2$. Jego grupa jest generowana przez automorfizm
Frobeniusa $\sigma(x)=x^2$. Ponieważ $\theta^2=\theta+1$,
mamy $\theta^3=\theta(\theta+1)=\theta^2+\theta=1$.
Grupa $L^\times$ składa się dokładnie z trzech elementów
$1,\theta,\theta^2$. Dla każdego $a\in L^\times$ zachodzi
$N_{L/K}(a)=a\sigma(a)=a^3=1$.

Sprawdźmy konstrukcję z 90. twierdzenia Hilberta dla
$a=\theta^2$. Wybieramy $b=\theta$. Wówczas
$\sigma(b)=\theta^2$ i
\[
\frac{b}{\sigma(b)}=\frac{\theta}{\theta^2}
=\theta^{-1}=\theta^2=a.
\]
W tym przykładzie jądro normy jest całą grupą $L^\times$.
Nie oznacza to, że norma każdego rozszerzenia ciał jest
identycznie równa $1$: zależy to od konkretnego rozszerzenia.
\end{przyklad}

\section{Addytywna postać 90. twierdzenia Hilberta}
\begin{twierdzenie}[90. twierdzenie Hilberta — postać addytywna]\label{aa:tw:hilbert90-additive}
Niech $L/K$ będzie skończonym cyklicznym rozszerzeniem Galois
stopnia $n\geq1$, a $\sigma$ generatorem $\operatorname{Gal}(L/K)$.
Dla każdego $a\in L$ następujące warunki są równoważne:
\begin{enumerate}[label=\textup{(\roman*)}]
\item $\operatorname{Tr}_{L/K}(a)=0$;
\item istnieje $b\in L$ takie, że $a=b-\sigma(b)$.
\end{enumerate}
\end{twierdzenie}
\begin{proof}
\emph{Kierunek (ii)$\Rightarrow$(i).}
Jeśli $a=b-\sigma(b)$, to z lematu
\ref{aa:lem:galois-norm-product} wynika
\[
\operatorname{Tr}_{L/K}(a)
=\sum_{j=0}^{n-1}(\sigma^j(b)-\sigma^{j+1}(b))
=b-\sigma^n(b)=0.
\]

\emph{Kierunek (i)$\Rightarrow$(ii).}
Rozpatrzmy $K$-liniowe odwzorowanie
$T:L\to L$, $T(b)=b-\sigma(b)$.
Jego jądro jest zbiorem elementów ustalonych przez $\sigma$,
czyli $L^{\langle\sigma\rangle}=K$ na mocy odpowiedniości Galois
(zob. twierdzenie~\ref{aa:tw:fundamental-galois}).
Ponieważ $[L:K]=n$, twierdzenie o rzędzie i jądrze daje
$\dim_K\operatorname{im}T=n-1$.

Ślad $\operatorname{Tr}_{L/K}:L\to K$ jest funkcjonałem
$K$-liniowym i jest niezerowy. Uzasadnijmy ten drugi fakt bez
założenia, że $\operatorname{char}K$ nie dzieli $n$:
niezdegenerowanie formy $L\times L\to K$,
$(x,y)\mapsto\operatorname{Tr}_{L/K}(xy)$,
wynikające z rozdziału~\ref{aa:ch:trace-norm}
(twierdzenie~\ref{aa:tw:trace-nondegenerate}), oznacza, że
istnieje $y\in L$ z $\operatorname{Tr}_{L/K}(1\cdot y)\neq0$.
Dlatego $\ker\operatorname{Tr}_{L/K}$ ma wymiar $n-1$.
Pierwszy kierunek dowodu wykazał już
$\operatorname{im}T\subseteq\ker\operatorname{Tr}_{L/K}$.
Obie są podprzestrzeniami tej samej przestrzeni $L$
o jednakowym skończonym wymiarze, a więc są równe.
Dla każdego $a$ o śladzie zerowym istnieje zatem $b\in L$
spełniające $T(b)=a$, zgodnie z tezą.
Dla $n=1$ obie podprzestrzenie są zerowe i dowód nadal działa.
\end{proof}

\begin{przyklad}[Addytywne równanie w charakterystyce dwa]
Niech $K=\mathbb F_2$ i
$L=\mathbb F_2[\theta]$, gdzie $\theta^2+\theta+1=0$.
Wielomian $X^2+X+1$ nie ma pierwiastka w $\mathbb F_2$,
więc $[L:K]=2$ i $L$ jest ciałem czteroelementowym.
Automorfizm $\sigma(x)=x^2$ ma rząd $2$ i ustala $K$;
w szczególności $\sigma(\theta)=\theta^2=\theta+1$.
Mamy $\operatorname{Tr}_{L/K}(1)=1+1=0$ oraz
$\operatorname{Tr}_{L/K}(\theta)=\theta+(\theta+1)=1$.
Stąd $\ker\operatorname{Tr}_{L/K}=\{0,1\}$, a
$\theta-\sigma(\theta)=\theta-(\theta+1)=1$,
ponieważ w charakterystyce $2$ zachodzi $-1=1$.
Zatem równanie $1=b-\sigma(b)$ ma rozwiązanie $b=\theta$,
mimo że $\operatorname{Tr}_{L/K}(1)=0$.
Przykład pokazuje, dlaczego w dowodzie nie wolno uzasadniać
niezerowości śladu wyłącznie równością
$\operatorname{Tr}_{L/K}(1)=[L:K]$.
\end{przyklad}

\begin{uwaga}[Dlaczego potrzebna jest cykliczność?]\label{aa:uw:hilbert90-cyclic}
Jeżeli skończone rozszerzenie Galois $L/K$ ma grupę
$G$, która nie jest cykliczna, żaden pojedynczy element
$\sigma\in G$ nie generuje całej grupy. Wtedy ciało stałe
$L^{\langle\sigma\rangle}$ jest większe od $K$.
Ze wzoru na wymiar obrazu operatora $1-\sigma$ wynika
$\dim_K\operatorname{im}(1-\sigma)<[L:K]-1$,
podczas gdy jądro niezerowego funkcjonału śladu ma wymiar
$[L:K]-1$. Nie można więc zamienić grupy cyklicznej
w twierdzeniu addytywnym na dowolną grupę Galois,
zachowując jedną ustaloną różnicę $b-\sigma(b)$.
\end{uwaga}

\begin{przyklad}[Kontrprzykład dla jednej różnicy w grupie niecyklicznej]\label{aa:prz:hilbert90-noncyclic}
Weźmy $L=\mathbb Q(\sqrt2,\sqrt3)$ oraz $K=\mathbb Q$.
Elementy $1,\sqrt2,\sqrt3,\sqrt6$ tworzą bazę $L$ nad $K$,
a grupa Galois jest izomorficzna z $C_2\times C_2$.
Niech $\sigma$ zmienia znak $\sqrt2$ i ustala $\sqrt3$.
Dla dowolnego
$x=x_0+x_1\sqrt2+x_2\sqrt3+x_3\sqrt6$,
$x_j\in\mathbb Q$, mamy
\[
(1-\sigma)(x)=2x_1\sqrt2+2x_3\sqrt6.
\]
Stąd $\operatorname{im}(1-\sigma)$ jest dwuwymiarową
podprzestrzenią rozpiętą przez $\sqrt2,\sqrt6$.
Z kolei suma czterech sprzężeń elementu $x$ wynosi $4x_0$,
więc $\ker\operatorname{Tr}_{L/K}$ jest trójwymiarową
podprzestrzenią rozpiętą przez $\sqrt2,\sqrt3,\sqrt6$.
W szczególności $a=\sqrt3$ ma ślad równy $0$,
ale nie istnieje $b\in L$ z $a=b-\sigma(b)$:
we wspomnianej bazie każda taka różnica ma współczynnik
przy $\sqrt3$ równy $0$.
Ten kontrprzykład nie przeczy twierdzeniu
\ref{aa:tw:hilbert90-additive}, ponieważ jego założenie
cykliczności grupy Galois nie jest spełnione.
\end{przyklad}